%%%%%%%%%%%%
%
% $Autor: OpenCode $
% $Pfad: Introduction.tex $
%
%%%%%%%%%%%%

\chapter{Introduction}
\label{ch:manual-introduction}

This manual explains how to operate the Germany Consumer Price Index (CPI) forecasting dashboard from a user perspective. The application is not a general-purpose data-science notebook or a developer toolkit. It is a focused Streamlit dashboard that lets the user inspect forecasts, compare forecasting models, adjust confidence intervals, and interpret the resulting graphs and evaluation metrics without editing Python source code.

The manual is written for three practical situations. First, it helps a first-time user launch the dashboard and confirm that the system is working correctly. Second, it guides a normal operator through the most important interface elements such as model selection, comparison mode, metric interpretation, and chart interaction. Third, it provides support-oriented sections for troubleshooting, operational safety, and technical reference so that the dashboard can be used confidently during demonstrations, project review, and day-to-day study work.

\begin{tcolorbox}[colback=blue!5!white, colframe=blue!75!black, title=Welcome to the CPI Forecasting Dashboard Manual]
This manual helps you start the dashboard, understand the controls, interpret the graph and metrics, and solve the most common user-side problems.

\textbf{Key user goals:}
\begin{itemize}
    \item launch the dashboard correctly,
    \item compare ARIMA, ETS, and SARIMA views,
    \item adjust confidence intervals,
    \item use chart interactions such as legend selection and plot download.
\end{itemize}
\end{tcolorbox}

The chapter order follows the normal user journey. The next chapter provides a quick-start path for users who want the shortest route to a successful first launch. After that, the manual explains the main function of the system, the installation and first-launch procedure, the dashboard interface, the available functions, help and support channels, troubleshooting, safety guidelines, technical specifications, and finally maintainer-oriented notes and optional extensions.
